\subsubsection{Circuits see only colors}\label{ov:types}

\paragraph{Color-determined relations.}
A relation $X\subseteq V^2$ is
\defn{color-determined} if
$\chi(u,v)=\chi(u',v') \implies X(u,v)=X(u',v')$ holds.
Here loops carry the separate label $\id$.
Equivalently, $X$ is a union of edge-color
classes, possibly together with the diagonal.
Write $X[\mu]$ for its common value on edges
of color $\mu$, $X[\id]$ for its value on
loops.

Every gate of an arbitrary \CoR\ circuit computes a color-determined relation on $\mathcal M_0$ (Lemma~\ref{lem:tables}).
We can prove this by induction on the
circuit gates. Input relations and constants are
color-determined, while Boolean operations and
converse preserve this property.
The only nontrivial step is composition $X;Y$.

For distinct $u,v$, the two-witness property realizes every pair of colors $\mu,\nu$ along a path $u\xrightarrow{\mu}z\xrightarrow{\nu}v$ with $z\notin\set{u,v}$.
%
The case $z\in\set{u,v}$ depends only on $\chi(u,v)$ and the diagonal values of $X,Y$.
Hence $(X;Y)(u,v)$ depends only on $\chi(u,v)$.
For a loop $(u,u)$, both edges of a path $u-z-u$ have the same color, and every color occurs at every vertex. 
Together with the case $z=u$, this makes $(X;Y)(u,u)$ independent of $u$, completing the induction.


\paragraph{Predicting and noticing a fold.}
A color-determined relation $X$ assigns a value $X[\mu]$ to each color $\mu$.
To predict the effect of a fold, we keep these values fixed and read them at the new edge colors:
\[
\defn{(Prediction)}\quad
X^{(\lambda)}(u,v):=X[\chi_\lambda(u,v)].
\]
Thus a recolored edge reads $X[\lambda^c]$ rather than $X[\lambda]$; all other pairs retain their values.
%
This predicts what a gate would compute if its color-dependent behavior on $\mathcal M_0$ remained unchanged after the fold;
its actual value on $\mathcal M_\lambda$ may differ.

Input relations and constants behave exactly as predicted.
Moreover, Boolean operations and converse commute with prediction
(Proposition~\ref{prop:prediction-commutation}); in particular,
$\overline{X^{(\lambda)}}=(\overline X)^{(\lambda)}$.
Unlike the preservation under embeddings~\eqref{eq:mono} in \S\ref{sec:natural-approaches},
this equality survives unrestricted complement.

Figure~\ref{fig:prediction} compares three operations for any two
distinct colors $\lambda,\mu$: $X$ selects the $\lambda$-colored
edges, and $Y$ the $\mu$-colored edges (both exclude loops).
The thick edge $e=(\hub,v_\mu)$ is unchanged by the fold of
$\lambda$, which removes all $X^{(\lambda)}$-edges at $\hub$.\nakamura{Fig.~\ref{fig:prediction} の表が読めない気がします。列見出しをもう少しちゃんと書くか、caption か本文に説明が欲しいです。
また、Fig.~\ref{fig:prediction} は Fig.~\ref{fig:fold} とほぼ同じにもみえます（あと、塗り替えた辺が Fig.~\ref{fig:prediction} では破線、Fig.~\ref{fig:fold} では実線で、整合性が取れていない）}

\begin{figure}[H]
  \centering
  \begin{minipage}[c]{0.58\linewidth}
    \centering
\begin{tikzpicture}[
      >=Latex,
      vertex/.style={circle,draw,fill=white,inner sep=0pt,minimum size=3.2pt},
      edge/.style={line width=.5pt},
      read/.style={line width=1.05pt},
      folded/.style={draw=red!65!black,densely dashed,line width=.75pt},
      every node/.style={font=\scriptsize},
    ]
      % These four vertices show two of the witnesses; the full arena
      % contains further vertices and satisfies the two-witness property.
      \coordinate (h0) at (0,0);
      \coordinate (v0) at (2.65,0);
      \coordinate (z0) at (1.325,0.98);
      \coordinate (w0) at (1.325,-0.98);
      \draw[read] (h0)--node[below,pos=.54] {$\mu$} (v0);
      \draw[edge] (h0)--node[above left] {$\lambda$} (z0);
      \draw[edge] (h0)--node[below left] {$\lambda$} (w0);
      \draw[edge] (z0)--node[above right] {$\mu$} (v0);
      \draw[edge] (w0)--node[below right] {$\mu$} (v0);
      \node[vertex,label=left:{$\hub$}] at (h0) {};
      \node[vertex,label=right:{$v_\mu$}] at (v0) {};
      \node[vertex,label=above:{$z$}] at (z0) {};
      \node[vertex,label=below:{$z'$}] at (w0) {};
      \node[font=\small] at (.1, 1.2) {$\mathcal M_0$};
      \draw[-{Latex[length=2mm]},line width=.7pt]
        (3.4,0)--node[above] {fold $\lambda$} (4.85,0);
      \begin{scope}[xshift=5.30cm]
        \coordinate (h1) at (0,0);
        \coordinate (v1) at (2.65,0);
        \coordinate (z1) at (1.325,0.98);
        \coordinate (w1) at (1.325,-0.98);
        \draw[read] (h1)--node[below,pos=.54] {$\mu$} (v1);
        \draw[folded] (h1)--node[above left,black] {$\lambda^c$} (z1);
        \draw[folded] (h1)--node[below left,black] {$\lambda^c$} (w1);
        \draw[edge] (z1)--node[above right] {$\mu$} (v1);
        \draw[edge] (w1)--node[below right] {$\mu$} (v1);
        \node[vertex,label=left:{$\hub$}] at (h1) {};
        \node[vertex,label=right:{$v_\mu$}] at (v1) {};
        \node[vertex,label=above:{$z$}] at (z1) {};
        \node[vertex,label=below:{$z'$}] at (w1) {};
        \node[font=\small] at (.1, 1.2) {$\mathcal M_\lambda$};
      \end{scope}
    \end{tikzpicture}
  \end{minipage}\hfill%
  \begin{minipage}[c]{0.40\linewidth}
    \centering
    {\footnotesize
    \renewcommand{\arraystretch}{1.16}%
    \begin{tabular}{@{}lccc@{}}
      \toprule
      Operation & Predicted & Actual & Notice? \\
      \midrule
      $X;Y$         & $1$ & $0$ & Yes \\
      $X\cup Y$     & $1$ & $1$ & No \\
      $\overline X$ & $1$ & $1$ & No \\
      \bottomrule
    \end{tabular}}
  \end{minipage}
  \caption{At $e=(\hub,v_\mu)$, only composition notices the fold of $\lambda$.
    Only two witness paths are shown; all $\lambda$-edges at $\hub$ are recolored to $\lambda^c$.}
  \label{fig:prediction}
\end{figure}

\if0
An operation \defn{notices} a fold at given color-determined
operands if applying it after prediction differs from predicting its
result. In Figure~\ref{fig:prediction}, composition notices the fold:
$X^{(\lambda)};Y^{(\lambda)}\ne(X;Y)^{(\lambda)}$.
Union and complement do not notice it; indeed, Proposition~\ref{prop:prediction-commutation}
shows that no Boolean operation or converse can ever notice a fold.
\fi

An operation \defn{notices} a fold at given
color-determined operands if applying it to
their predictions differs from predicting
its result. In Figure~\ref{fig:prediction}, composition notices
the fold at $e$, since
\[
  e\notin\bigl(X^{(\lambda)};Y^{(\lambda)}\bigr)
  \quad \text{and} \quad
  e\in(X;Y)^{(\lambda)}.
\]
By contrast, union and complement do not
notice the fold. More generally, no Boolean
operation or converse can notice any fold
(Proposition~\ref{prop:prediction-commutation}).

Every fold must be noticed by some composition gate.
Let $C$ be a \CoR\ circuit such that $C\equiv G_k$.
If no composition gate notices the fold of $\lambda$,
then, since $(\hub,v_{\lambda^c})$ is unchanged,
Lemma~\ref{lem:unnoticed} gives
\[
  \mathcal M_0\models C(\hub,v_{\lambda^c})
  \quad\Longleftrightarrow\quad
  \mathcal M_\lambda\models C(\hub,v_{\lambda^c}).
\]
This contradicts $C\equiv G_k$, since
$\mathcal M_0\models G_k(\hub,v_{\lambda^c})$
but $\mathcal M_\lambda\not\models G_k(\hub,v_{\lambda^c})$.

\if0
Every fold must be noticed by some composition gate.
Now let $C$ be any circuit computing $G_k$.
If no composition gate notices the fold of
$\lambda$, Lemma~\ref{lem:unnoticed} implies that the output
on $\mathcal M_\lambda$ is the prediction
of the output on $\mathcal M_0$.
Since the test edge $e_\lambda=(\hub,v_{\lambda^c})$ is unchanged,
\[
\sem{C}{\mathcal M_\lambda}(e_\lambda)
=
\bigl(\sem{C}{\mathcal M_0}
\bigr)^{(\lambda)}(e_\lambda)
=
\sem{C}{\mathcal M_0}(e_\lambda),
\]
contradicting the test for $\lambda$, whose outputs are $0$ on $\mathcal M_{\lambda}$ and $1$ on $\mathcal M_0$.
\fi
